%% ============================================================
%% 分片 33：原文第 11 节「$A$-Branes and ${\cal D}$-Modules」（原文 glshort.tex 起于第 9938 行）
%% 由 tools/harvmac2tex.py 从 arXiv 源码（hep-th/0604151v3）自动生成。
%% 只译英文散文；公式、编号、\label、\cite 一律不改。
%% ============================================================

\subsection{对应于 $A$-膜的 ${\cal D}$-模}
\label{sec:dmab}

c.c. 膜的存在意味着：复结构 $K$ 中的每个 $A$-膜都是 $K_\M^{1/2}$ 上微分
算子层的一个模。事实上，一般地，若 ${\cal B}$ 与 ${\cal B}'$ 是任意两个
膜，则 $({\cal B},{\cal B}')$ 弦构成 $({\cal B},{\cal B})$ 弦之代数的
一个模。这一想法在第~\ref{sec:comptwo}~节 的 \dunamp\ 中作过说明，
并且是开-闭拓扑场论通常公理的一部分 \cite{ref142,ref143}。在层化可行的
范围内，$({\cal B},{\cal B})$ 弦构成一个代数\emph{层}，而不只是一个
代数，且 $({\cal B},{\cal B}')$ 弦构成该代数层的一个模层。这一陈述的
意思只是：环结构与模结构可以对仅在靶空间中某个适当开集内正则的开弦
定义。

我们想把这一构造应用到 ${\cal B}={\cal B}_{c.c.}$ 为典范余各向同性膜的
情形。在此情形，我们断言：对任何在复结构 $K$ 中为 $A$-型的膜
${\cal B}$，$({\cal B}_{c.c.},{\cal B}')$ 弦都没有瞬子修正。这一论证
也适用于涉及 c.c. 膜的高阶拓扑耦合（如三次 Yukawa 耦合）。无瞬子修正
这一点可以按与完全落在 c.c. 膜上的弦大致相同的方式论证。相关的圆盘
瞬子有一部分边界落在 c.c. 膜上，而在这部分边界上瞬子必须是常数。但由
解析性，该瞬子必须处处为常数。

从先前的讨论我们还知道，何种层化是可行的。我们可以把一个代数与一个模
联系到 $\MH$ 中形如 $T^*{{\cal V}}$、${{\cal V}}\subset \M$ 的开集上，
但未必能联系到更一般的开集上。因此，复结构 $K$ 中 $\MH(G,C)$ 上的一个
$A$-膜给出 $\M$ 上代数层 ${\cal D}_{K_{
\EUM}^{1/2}}$ 的一个模层。

现在我们给出几个例子。（关于涉及余各向同性膜的开弦量子化的若干详细
例子，见 \cite{ref144}。）

我们的第一个例子是由条件 $\varphi=0$ 定义、且 Chan-Paton 丛平凡的
$A$-膜 ${\cal B}'$。该子流形是 $\cM(G,C)$ 的一个拷贝，在复结构 $K$ 与
$J$ 中是 Lagrange 子流形，在复结构 $I$ 中是复子流形。因此 ${\cal B}'$
是一个 $(B,A,A)$-膜的例子。（我们在第~\ref{sec:opprod}~节中分析
't Hooft 算子的算子乘积展开时用过这个膜。）为简单起见，我们取
${\cal B}'$ 上的平坦 Chan-Paton 联络为平凡联络。如式
\eqref{eq:exactcc} 所述，c.c. 膜有一个 Chan-Paton 线丛 ${\cal N}$，
它拓扑平凡，但带有非平凡的联络。

要计算经典层面的开弦谱，可以把超对称 sigma 模型约化为超对称量子力学，
即只保留零模。在这一近似下，具有 $({\cal B}_{c.c.},{\cal B}')$ 边界
条件的开弦态，是某个向量丛与 ${\cal N}^{-1}$ 的张量积的截面，该向量丛
由对费米零模空间作量子化得到。${\cal N}$ 上的联络限制到 $\varphi=0$
时为零，所以可以略去 ${\cal N}^{-1}$ 这个因子。费米零模空间就是
$\M\subset \MH$ 的切空间（与 $\M$ 的法向相伴的费米子在弦的两端满足
相反的边界条件，因而没有零模）。对费米零模空间作量子化，就得到 $\M$
上的旋量。把 $\M$ 看作复流形，它的旋量丛就是
$K_\M^{1/2}\otimes(\oplus_{i=0}^{{\rm dim}\,\M}\Omega^{0,j}(\M))$。
这里 $\Omega^{0,j}$ 是 $\M$ 上 $(0,j)$-形式的层。

在此情形，BRST 算子或拓扑超对称 $Q$ 就是作用在取值于 $K_\M^{1/2}$ 的
$(0,j)$-形式上的 $\bar\partial$ 算子。因此，物理的
$({\cal B}_{c.c.},{\cal B}')$ 弦的层，就是\footnote{$j>0$ 的
$\Omega^j$ 形式对足够小的开集 ${{\cal V}}$ 上的上同调没有贡献，故在此
可以略去。} $K_\M^{1/2}$ 的全纯截面层。这当然是环层
${\cal D}_{K_{ \EUM}^{1/2}}$ 的一个模层；事实上，这个环层正是通过它
（作为其模层）来定义的。这一构造也许更直接地解释了为什么由 c.c. 膜
导出的环层恰好是 ${\cal D}_{K_{ \EUM}^{1/2}}$，而不是对某个别的
${\cal L}$ 的 ${\cal D}_{\cal L}$。

我们声称费米零模空间的量子化恰好给出 $\M$ 上的旋量、而不是取值于某个
线丛 ${\cal L}$ 的旋量，这看起来也许又隐含地用到了时间反演对称性
$\EUT$。其实，对这一点我们可以给出另一个论证，尽管它用到该理论的另一
个离散对称性。在定义于 $M=\Bbb{R}\times I\times C$ 的四维超对称规范
理论中，膜 ${\cal B}_{c.c.}$ 与 ${\cal B}'$ 是物理上合理的酉膜；这里
$\Bbb{R}$ 是时间方向，$I$ 是一个区间，其两端的边界条件分别由
${\cal B}_{c.c.}$ 与 ${\cal B}'$ 定义。为在 $M$ 上定义该理论，需要
一些扭化以在 $C$ 上的紧化中保持超对称，但没有涉及时间方向的扭化。
所以除了拓扑超对称 $Q$ 之外，$M$ 上的理论还在伴随超对称 $Q^\dagger$
下不变（二者满足 $\{Q,Q^\dagger\}=H$，$H$ 为哈密顿量）。物理理论具有
一个「电荷共轭」对称性，它交换 $Q$ 与 $Q^\dagger$。在这一对称性下
不变，意味着对费米零模作量子化得到的 Chan-Paton 丛恰好是 $\CM$ 的
旋量丛，而不是旋量丛与某个额外线丛的张量积。\footnote{若 $\pi_1(G)$
非平凡，可以修改该构造：取膜 ${\cal B}'$ 上的 Chan-Paton 丛为一个平坦
线丛。于是我们得到 $\M$ 上取值于一个平坦线丛的旋量；作为特例，这给出
$\M$ 上所有可能的自旋结构。由于 $D$-膜与 $K$-理论的关系
\cite{ref102}，$\M$ 的自旋结构的选取实际上是膜 ${\cal B}'$ 的
Chan-Paton 丛的细致描述的一部分。} 这类论证最早出现在
\cite{ref145}，用于量子化孤立子。

注意，在这一计算中重要的是：${\cal B}'$ 的支撑不仅是关于 $\omega_K$
的 Lagrange 子流形，而且还是关于由 c.c. 膜确定的复结构 $\hat N=I$ 的复
子流形。否则，拓扑超荷或 BRST 算子就不会化为复结构 $I$ 中的
$\bar\partial$ 算子。

我们的第二个例子是对几何 Langlands 纲领重要的情形：一个 $\CMF$ 型膜，
即支撑在 Hitchin 纤维化的一条纤维 $\CMF$ 上、带平坦 Chan-Paton 线丛的
膜 ${\cal B}_{\CMF}$。这是一个 $(A,B,A)$-型膜，所以特别地，它是复结构
$K$ 中的 $A$-膜。因此，它给出 ${\cal D}_{K_{ \EUM}^{1/2}}$ 的一个模层。
我们在第~\ref{sec:genotics}~节 末尾将通过进一步的基本论证，说明如何把
它变成 $\cal D$（即 $\M$ 上通常微分算子层）的一个模层。所以
$\CMF$ 型膜具有两个关键性质：它是磁特征膜，因为它是零膜的 $S$-对偶；
并且它给出 $\M$ 上 ${\cal D}$-模的层。这些正是几何 Langlands 纲领的
基本主张。

显式描述来自 $\CMF$ 型膜的 ${\cal D}$-模层或
${\cal D}_{K_{ \EUM}^{1/2}}$-模层是困难的。但在交换情形 $G=U(1)$ 中
我们可以做到这一点。在这一特殊情形，sigma 模型是自由场论，且边界条件
是线性的，所以开弦的谱与模结构的计算是直接的。

对 $G=U(1)$，我们可以把 $A$ 与 $\phi$ 看作实一形式（当然 $A$ 只在
规范变换下定义，而 $\phi$ 是规范不变的）。$A$ 与 $\phi$ 的 Hitchin
方程解耦，而在复结构 $I$ 中，Hitchin 模空间是 $C$ 的 Jacobi 簇
${\rm Jac}(C)$（即 $C$ 上拓扑平凡的全纯线丛的模空间）与向量空间
${\EUBB}=H^0(C,\Omega^1)$ 的乘积。${\rm Jac}(C)$ 是维数为 $g_C$
（$C$ 的亏格）的复环面，而 Hitchin 模空间
${\rm Jac}(C)\times {\EUBB}$ 可以等同于它的余切丛。Hitchin 纤维化就是
到 ${\EUBB}$ 的投影。

给定 $p\in {\EUBB}$ 以及 Hitchin 纤维化相应的纤维
${\CMF}_p\simeq {\rm Jac}(C)$，我们可以像上一个例子那样计算
$({\cal B}_{c.c.},{\cal B}_{{\CMF}_p})$ 弦的空间。一个差别是：虽然膜
${\cal B}_{{\CMF}_p}$ 的 Chan-Paton 线丛 $\cL\to{\CMF}_p$ 仍然拓扑
平凡，但我们现在允许它带有任意的酉平坦联络。另一个差别是：c.c. 膜的
Chan-Paton 线丛 ${\cal N}$ 限制到 ${\CMF}_p$ 上时，一般是全纯非平凡的。
另一方面，由于 ${\CMF}_p$ 是环面，其典范丛 $K_{{\CMF}_p}$ 是平凡的。
因此，鬼数为零的 BRST 不变开弦态是
$T_p=\cL\otimes {\cal N}^{-1}\vert_{{\CMF}_p}$ 的全纯截面。

但是，$({\cal B}_{c.c.},{\cal B}_{c.c.})$ 弦的层是如何作用在 $T_p$ 的
截面上的呢？我们下面回答这个问题。我们在 $\RR\times I$ 上工作，左边界
落在膜 ${\cal B}_{c.c.}$ 上，右边界落在 ${\cal B}_{{\CMF}_p}$ 上。
Hitchin 纤维化的纤维 $p$ 由 $\phi=\upsilon$ 定义，其中
$\upsilon=\upsilon_z\,dz+\upsilon_\bz\,d\bz$ 是一个实调和一形式。膜
${\cal B}_{{\CMF}_p}$ 的 Chan-Paton 线丛 ${\cal L}\to{\CMF}_p$ 有一个
平坦酉联络，它可以方便地表示为 % 式 (11.17)
\setcounter{equation}{16}
\begin{equation}\label{eq:zoncob}
\beta={1\over
{2\pi}}\int_C |d^2 z|\left(a_z \delta A_\bz+a_\bz \delta
A_z\right)
\end{equation}
 其中 $a=a_z\,dz+a_\bz\,d\bz$ 是 $C$ 上的实调和一形式。类似地，c.c. 膜
上线丛 ${\cal N}$ 的联络是\footnote{为得到这一公式及下一个公式，我们
使用式~\eqref{eq:tomog} 中关于 $\omega_J$ 与 $\omega_K$ 的表达式。在
式~\eqref{eq:tomog} 中，${\rm Tr}$ 表示 $U(N)$ 的 $N$ 维表示中的迹，
且 $A$ 与 $\phi$ 被理解为反厄米的。由于这里我们把 $A$ 与 $\phi$ 看作
实一形式，必须补上一个负号。} % 式 (11.18)
\setcounter{equation}{17}
\begin{equation}\label{eq:homifly}
\alpha=-{\Im\,\tau\over
2\pi}\int_C |d^2 z|\left(\phi_z \delta A_\bz+\phi_\bz \delta
A_z\right).
\end{equation}
 最后，辛形式 $\omega=({\Im}\,\tau)\omega_K$ 是如下一形式的外微分 % 式 (11.19)
\setcounter{equation}{18}
\begin{equation}\label{eq:nomigly}
\zeta=-{i\,
\Im\,\tau\over 2\pi}\int_C |d^2 z|\left(\phi_\bz \delta A_z-\phi_z
\delta A_\bz\right).
\end{equation}

$A$-模型的作用量在模去 $Q$-恰当项之后，是 $\Sigma$ 边界上的积分。
作用量中来自右边界（我们称之为 $\partial_R\Sigma$）的贡献是 % 式 (11.20)
\setcounter{equation}{19}
\begin{equation}\label{eq:omitro}
\int_{\partial_R\Sigma}
\Phi^*(\zeta -i\beta)={i\,\Im\,\tau\over 2\pi}\int_{\partial_R\Sigma}
\left(\left(\upsilon_z-(\Im\,\tau)^{-1}a_z\right) \dot{A}_\bz
-\left(\upsilon_\bz+(\Im\,\tau)^{-1}a_\bz\right) \dot{A}_z\right)ds.
\end{equation}

这里我们用了右边界上 $\phi=\upsilon$ 这一事实。此外，$s$ 是
$\Sigma=\Bbb{R}\times I$ 上的「时间」坐标，$\dot A=\partial A/\partial s$。
左边界 $\partial_L\Sigma$ 的定向相反，其贡献为
% 式 (11.21)
\setcounter{equation}{20}
\begin{equation}\label{eq:bofloc}
-\int_{\partial_L\Sigma} \Phi^*(\zeta-i\alpha)=-{i\,\Im\,\tau\over \pi}\int_{\partial_L\Sigma}
 \phi_z
\dot{A}_\bz\,ds.
\end{equation}

超对称弦态来自沿弦方向的零模，所以为确定它们，$\partial_L\Sigma$ 与
$\partial_R\Sigma$ 的区分并不重要，我们只需把两项贡献加到作用量中。
总作用量因而为
% 式 (11.22)
\setcounter{equation}{21}
\begin{equation}\label{eq:nofloc}
{i\,\Im\,\tau\over \pi}\int \left(-\phi_z+{1\over
2}(\upsilon_z-(\Im\,\tau)^{-1}a_z)\right)\dot{A}_\bz \,ds-{i\,\Im\,\tau\over
2\pi}\int(\upsilon_\bz+(\Im\,\tau)^{-1}a_\bz) \dot{A}_z\,ds.
\end{equation}

我们看到，作用量依赖参数 $\upsilon_z,\upsilon_\bz,a_z,a_\bz$ 的组合
% 式 (11.23)
\setcounter{equation}{22}
\begin{equation}\label{eq:numiko}
\upsilon_z-(\Im\,\tau)^{-1}a_z,\ \upsilon_\bz+(\Im\,\tau)^{-1}a_\bz.
\end{equation}

式~\eqref{eq:nofloc} 中的第一项在量子化后告诉我们 $\phi_z$ 如何作用：
% 式 (11.24)
\setcounter{equation}{23}
\begin{equation}\label{eq:xococox}
\phi_z\to -i\pi(\Im\,\tau)^{-1}\left({\delta\over\delta
A_\bz}+{i\,\Im\,\tau\over 2\pi}\upsilon_z-{i\over 2\pi} a_z\right).
\end{equation}

这是 ${\rm Jac}(C)$ 上某个拓扑平凡线丛上的协变 $\partial$ 算子。第二项
表明波函数是一个拓扑平凡但全纯非平凡的线丛
$T_p\to{\rm Jac}(C)$ 的截面，这个线丛实际上就是我们先前认出的同一个
线丛 $T_p=\cL\otimes {\cal N}^{-1}\vert_{{\CMF}_p}$。事实上，这一项是
全导数，所以若把它吸收到初态与末态的波函数中，我们就会发现波函数不再
与 $A_z$ 无关，而是被如下算子零化
% 式 (11.25)
\setcounter{equation}{24}
\begin{equation}\label{eq:bilbox}
{\delta\over\delta A_z}-{i\,\Im\,\tau\over 2\pi}\upsilon_\bz-{i\over
2\pi} a_\bz .
\end{equation}

这是 ${\rm Jac}(C)$ 上平凡线丛上的协变 $\bpartial$ 算子。显然，这两个
算子定义了 ${\rm Jac}(C)$ 上的一个平坦联络。它酉当且仅当 $\upsilon=0$。

总之，对 $({\cal B}_{c.c.},{\cal B}_{{\CMF}_p})$ 弦作量子化，给了我们
一个 ${\cal D}$-模，它对应于 ${\rm Jac}(C)$ 上平凡线丛的一个复平坦联络
的选取。另一方面，$S$-对偶把这一族 $A$-膜等同于 $\MH(U(1),C)$ 上所有
零膜的集合，而 $\MH(U(1),C)$ 在复结构 $J$ 下是 $C$ 上平坦
$\Bbb{C}^*$-丛的模空间。换言之，在交换情形，$S$-对偶在 $C$ 上平坦
$\CC^*$ 联络的规范等价类与 ${\rm Jac}(C)$ 上平坦 $\CC^*$ 联络的规范
等价类之间建立了一个自然对应。

这一对应可以直接看出。流形 $M$ 上的平坦 $\CC^*$ 联络等同于从
$H_1(M,\Bbb{Z})$ 到 $\CC^*$ 的一个同态。因此只需证明
$H_1(C,\Bbb{Z})$ 同构于 $H_1({\rm Jac}(C),\Bbb{Z})$。但由于
${\rm Jac}(C)\simeq H^1(C,\RR)/H^1(C,\Bbb{Z})$，这等同于证明
$H_1(C,\Bbb{Z})\simeq H^1(C,\Bbb{Z})$，而这就是 Poincaré 对偶。

\begin{yizhu}[易错点]
按：第一段要分清「环层」与「模层」两个方向。c.c. 膜对自身给出的是
\emph{环}层 ${\cal D}_{K_\M^{1/2}}$，即 $({\cal B}_{c.c.},
{\cal B}_{c.c.})$ 弦；而任意 $K$ 中 $A$-型膜 ${\cal B}$（或
${\cal B}'$）给出的是该环层上的\emph{模}层，即
$({\cal B}_{c.c.},{\cal B})$ 弦。原文「每个 $A$-膜都是微分算子层的一个
模」说的正是这个模层方向，切勿反过来说「c.c. 膜是某膜的模」。此外，
层化只在 $T^*{{\cal V}}\subset\MH$（${{\cal V}}\subset\M$）型开集上
可行，这一点承接前一节的部分层化，并不是全局层化。
\end{yizhu}

\begin{yizhu}[易错点]
按：式 \eqref{eq:homifly} 的负号来自约定而非笔误，脚注已作说明：前文
式~\eqref{eq:tomog} 把 $A,\phi$ 取为反厄米、${\rm Tr}$ 取 $N$ 维表示
中的迹；本片为方便把 $A,\phi$ 当作实一形式，所以要补一个负号。注意
式 \eqref{eq:zoncob} 的 $\beta$ 就没有负号（它来自 $C$ 上的实调和
一形式 $a$），而式 \eqref{eq:homifly} 的 $\alpha$ 带负号，与前一节的
式~\eqref{eq:exactcc}（32 片中的 (11.3)）恰好差一个符号。三式对照着看
可避免把符号当成排版错误。
\end{yizhu}

\begin{yizhu}[原文疑误]
按：旋量丛的分解
$K_\M^{1/2}\otimes(\oplus_{i=0}^{{\rm dim}\,\M}\Omega^{0,j}(\M))$ 中，
求和指标写作 $i$ 而项中却用 $j$，二者不一致；按上下文（$\Omega^{0,j}$
是 $(0,j)$-形式的层，$j$ 从 $0$ 到 ${\rm dim}\,\M$），求和应为
$\oplus_{j=0}^{{\rm dim}\,\M}\Omega^{0,j}(\M)$。译文照录原文，请读者
按 $j$ 理解。此式给的是复流形 $\M$ 的旋量丛：$\M$ 是凯勒流形，
$K_\M^{1/2}$ 配上 $(0,j)$-形式正是旋量丛的分解。
\end{yizhu}

\begin{yizhu}[原文疑误]
按：本片另有两处英文笔误。其一，交换情形一段写作「in the the abelian
case」，冠词重复；其二，回答 $({\cal B}_{c.c.},{\cal B}_{c.c.})$ 弦
如何作用时写作「the sheaf of ... strings acts」（主语为 sheaf，却用
复数动词 acts）。译文按正确语法处理，特此说明，以免读者以为原文另有所
指。
\end{yizhu}

\begin{yizhu}[含糊点拨]
按：何以外围费米子没有零模？因为 $\M\subset\MH$ 的切向费米子在弦的
两端满足同一边界条件，可以有无零模的常值模式；但法向费米子（对应于
$\M$ 在 $\MH$ 中的法向）在两端满足相反的边界条件，只能取零。于是费米
零模空间就是切空间 $T\M$，量子化它得到 $\M$ 上的旋量。还有一个隐含
要点：$Q$ 要在这一步化为复结构 $I$ 中的 $\bar\partial$，必须让 ${\cal
B}'$ 的支撑 $\M$ 关于 $\hat N=I$ 是复子流形（同时它已是关于
$\omega_K$ 的 Lagrange 子流形）。这正是原文强调「otherwise ... would
not reduce to the $\bar\partial$ operator in complex structure $I$」
的原因，也是 ${\cal B}'$ 为 $(B,A,A)$-膜这一事实的另一面。
\end{yizhu}

\begin{yizhu}[物理图像]
按：第二个例子是本文的主线：$\CMF$ 型膜支撑在 Hitchin 纤维化的一条
纤维 $\CMF$（紧致复环面）上、带平坦 Chan-Paton 线丛。它有双重身份：
一方面它是零膜的 $S$-对偶，即磁特征膜；另一方面它给出 $\M$ 上
${\cal D}$-模的层。原文说这两点「正是几何 Langlands 纲领的基本主张」，
因为几何 Langlands 的核心对应是：$C$ 上平坦 ${}^L G$-丛（局部系统）
对应于 $G$-丛的模空间 $\M(G,C)$ 上具 Hecke 特征的 ${\cal D}$-模。注意 $\CMF$ 是
Hitchin 纤维（粗体 $\cmmib F$），不要与 $\MH$、$\M$ 或膜符号
${\cal B}$ 混淆；参见 Frenkel 2009 综述中关于本文与几何 Langlands 对应的
讨论。
\end{yizhu}

\begin{yizhu}[物理图像]
按：交换情形的对应链条值得记住：$({\cal B}_{c.c.},
{\cal B}_{{\CMF}_p})$ 弦的量子化给出 ${\rm Jac}(C)$ 上平凡线丛的复
平坦联络所对应的 ${\cal D}$-模；$S$-对偶把它等同于 $\MH(U(1),C)$ 上
全部零膜的集合，而在复结构 $J$ 下 $\MH(U(1),C)$ 就是 $C$ 上平坦
$\CC^*$-丛的模空间。于是 $S$-对偶在「$C$ 上的平坦 $\CC^*$ 联络」与
「${\rm Jac}(C)$ 上的平坦 $\CC^*$ 联络」之间建立自然对应。最后一段
把这个对应化归为同调层面的恒等式：平坦 $\CC^*$ 联络等同于
$H_1(M,\Bbb{Z})\to\CC^*$ 的同态，而
${\rm Jac}(C)\simeq H^1(C,\RR)/H^1(C,\Bbb{Z})$，故只需
$H_1(C,\Bbb{Z})\simeq H_1({\rm Jac}(C),\Bbb{Z})$，即 Poincaré 对偶
$H_1(C,\Bbb{Z})\simeq H^1(C,\Bbb{Z})$。这正是「$S$-对偶归结为 $T$-对偶」
在交换情形的最简单体现。另外注意式~\eqref{eq:xococox} 与
\eqref{eq:bilbox} 给出的协变 $\partial$、$\bar\partial$ 算子定义了
${\rm Jac}(C)$ 上的平坦联络，它仅在 $\upsilon=0$ 时才是酉的，这也是
把这一族 $A$-膜与零膜对应的一个技术要点。
\end{yizhu}
